\documentclass[fleqn]{article}
\oddsidemargin 0.0in
\textwidth 6.0in
\thispagestyle{empty}
\usepackage{import}
\usepackage{amsmath}
\usepackage{graphicx}
\usepackage{flexisym}
\usepackage{calligra}
\usepackage{amssymb}
\usepackage{bigints} 
\usepackage[english]{babel}
\usepackage[utf8x]{inputenc}
\usepackage{float}
\usepackage[colorinlistoftodos]{todonotes}


\DeclareMathAlphabet{\mathcalligra}{T1}{calligra}{m}{n}
\DeclareFontShape{T1}{calligra}{m}{n}{<->s*[2.2]callig15}{}
\newcommand{\scriptr}{\mathcalligra{r}\,}
\newcommand{\boldscriptr}{\pmb{\mathcalligra{r}}\,}

\definecolor{hwColor}{HTML}{1a0252}

\begin{document}

  \begin{titlepage}

    \newcommand{\HRule}{\rule{\linewidth}{0.5mm}}

    \center

    \begin{center}
      \includegraphics[height=11cm, width=11cm]{asu.png}
    \end{center}

    \vline

    \textsc{\LARGE Quantum Physics II}\\[1.5cm]

    \HRule \\[0.5cm]
    { \huge \bfseries Problem Set Three}\\[0.4cm] 
    \HRule \\[1.0cm]

    \textbf{Behnam Amiri}

    \bigbreak

    \textbf{Prof: Onur Erten}

    \bigbreak

    \textbf{{\large \today}\\[2cm]}

    \vfill

  \end{titlepage}

  \begin{enumerate}
    \item \textbf{French problem 3-8} 
    Two inertial coordinate system $S$ and $S^'$ move with speed $c/2$ with 
    respect to each other. Draw a Minkowski diagram relating these two 
    systems. (Let the axes of $x$ and $ct$ for $S$ be at right angles in your drawing.)
    \begin{enumerate}
      \item Draw calibration hyperbolas that allow you to define unit distances along the axes of $x, x^', ct$ and $ct^'$.

      \item Plot the folowing point events on the diagram: $(1) ~ x=1, ct=1; ~ (2) ~ x^'=1, ct^'=1; ~ (3) ~ x^'=2, ct^'=0; ~ (4) ~ x=0, ct=2$.


      \item From your diagram determine the coordinates in $S^'$ (or $S$) corresponding to the above events.
      
    \end{enumerate}


    \item \textbf{French problem 5-20} $A$ and $B$ are twins. $A$ goes on a trip to $\alpha$ Centauri (4 light-years away) and back again.
    He travels at speed $0.6c$ with respect to the earth both ways, and transmits a radio signal every $0.01$ year in his frame.
    His twin $B$ similarly sends a signal every $0.01$ year in his own rest frame. 

    \begin{enumerate}
      \item How many signals emitted by $A$ before he turns around does $B$ receive?
      
      \item How many signals does $A$ receive before he turns around?
      
      \item What is the total number of signals each twin receives from other one?
      
      \item Who is younger at the end of the trip, and by how much? 
    \end{enumerate}
    Show that the twins both agree on this result.
    

    \item \textbf{French problem 6-4 part (a)} A photon of energy $E$ collides with a stationary particle of rest mass $m_0$ and 
    is absorbed. What is the velocity of the resulting composite particle?

      \textcolor{hwColor}{
        \\
        \textbf{Before the collision:}
        \\
        \\
        $
          \begin{cases}
            E_0=m_0 ~ c^2
            \\
            \\
            p_0=0
          \end{cases}
        $
        \\
        \\
        \\
        \textbf{After the collision:}
        \\
        \\
        $
          E_1=\gamma ~ M ~ c^2
          \\
          \\
          p_1=\gamma ~ M ~ c ~ \beta
        $
        \\
        \\
        Coservation of energy and momentum tells us:
        \\
        \\
        $
          \begin{cases}
            E+E_0=E_1
            \\
            p+p_0=p_1
          \end{cases} \Longrightarrow \begin{cases}
            E+m_0 ~ c^2=\gamma ~ M ~ c^2
            \\
            \dfrac{E}{c}+0=\gamma ~ M ~ c ~ \beta
          \end{cases}
           \Longrightarrow \begin{cases}
             \dfrac{\dfrac{E}{c}}{E+m_0 ~ c^2}=\dfrac{\gamma ~ M ~ c ~ \beta}{\gamma ~ M ~ c^2}
           \end{cases}
           \\
           \\
           \\
           \therefore ~~~ \boxed{\beta=\dfrac{1}{1+m_0 \dfrac{c^2}{E}}} ~~~~ \checkmark
        $
      }
    

  \end{enumerate}

\end{document}
